How short can a message be made? How fast can it travel through a noisy wire without a single error? Can a secret be
kept from someone who can compute anything? Shannon, Hamming and Diffie and Hellman turned these questions into
theorems, and this book proves them.

Along the way: entropy and the inequalities that govern it, compression and Kolmogorov complexity, the capacity of
noisy channels, Hamming and Golay codes and the lattice $E_8$ drawn on the cover, Reed--Solomon codes and list
decoding, sparse-graph codes and polar codes, one-way functions and pseudorandomness, RSA, elliptic curves and the
attacks on them, lattices and learning with errors, zero-knowledge proofs and secret sharing, and quantum computing
up to the algorithms of Shor and Grover.

Every theorem is proved in full or clearly labelled as quoted, every exercise has a hint and a reasoned solution, and
every chapter comes with a Python notebook that builds the coders, decoders and ciphers and runs them. Written for the
pleasure of seeing how far mathematics can carry a message.
